\documentclass{scdpg}
\begin{document}
\scBookLanguage{de}
\begin{scAbstract}
\scLanguage{en}
\scTitle{Atom-Photon Quantum Gate using an Atomic Clock Qubit}
\scAuthor{}{Leart}{Zuka}{1}
\scAuthor{}{Tobias}{Frank}{1}
\scAuthor{}{Gianvito}{Chiarella}{1}
\scAuthor{}{Pau}{Farrera}{1,2}
\scAuthor{}{Gehard}{Rempe}{1}
\scAffiliation{1}{Max Planck Institute for Quantum Optics, Hans-Kopfermann-Straße 1, 85748 Garching bei München}
\scAffiliation{2}{Munich Center for Quantum Science and Technology (MCQST), Schellingstr. 4, D-80799 München}
\scBeginText
Single atoms coupled to optical cavities provide a powerful platform for photonic quantum information processing. Minimizing the influence of external factors like magnetic field fluctuations is a key goal in this endeavor. In this work, we investigate a novel gate protocol that employs a single Rb$^{87}$ atom prepared in a magnetic-field-insensitive clock qubit and coupled to a high-finesse birefringent Fabry–Perot cavity. The cavity ’ s high birefringence creates two well-separated polarization eigenmodes, enabling polarization-selective reflectivity that depends on the atomic qubit state. This allows us to implement a CPHASE gate between the atomic clock states and our photonic qubits. We present simulations based on input–output theory and master-equation modeling that quantify the conditional reflection amplitudes, gate truth table, and resulting fidelities under realistic experimental imperfections such as finite mode matching, cavity loss channels, and multiphoton contributions. We further report on ongoing experimental progress towards implementing the protocol. Our results indicate that atom-state-dependent phase shifts on the photonic qubit are achievable in the current system, providing a viable path toward a robust, high-fidelity, cavity-assisted atom–photon quantum gate.
\scEndText
\scConference{Mainz 2026}
\scPart{Q}
\scContributionType{Poster}
\scTopic{4.4 Quantum Technologies - Networks, Repeaters, and QKD}
\scKeywords{Cavity QED; Fiber Cavities; Quantum Gate; Atom-Photon Gate}
\scEmail{leart.zuka@mpq.mpg.de}
\scCountry{Germany}
\end{scAbstract}
\end{document}
